% !TeX program = xelatex
% A canonical shear preserves phase-space area and Poisson structure.
% Compile: xelatex canonical_transformation_tikz.tex
\documentclass[tikz,border=6pt]{standalone}
\usepackage{amsmath}
\usepackage[UTF8,fontset=fandol]{ctex}
\usepackage{fix-cm}
\usetikzlibrary{arrows.meta}

\begin{document}
\begin{tikzpicture}[>=Stealth,thick,every node/.style={font=\fontsize{16}{19}\selectfont}]
  % Old coordinates.
  \begin{scope}
    \draw[blue!70!black,->] (-.3,0) -- (4.8,0) node[right] {$q$};
    \draw[blue!70!black,->] (0,-.3) -- (0,4.6) node[above] {$p$};
    \foreach \x in {1,2,3,4} {\draw[black!25] (\x,0) -- (\x,4);}
    \foreach \y in {1,2,3,4} {\draw[black!25] (0,\y) -- (4,\y);}
    \fill[blue!18] (1,1) rectangle (2,2);
    \draw[blue!70!black] (1,1) rectangle (2,2);
    \node[blue!70!black] at (2.15,4.75) {旧变量};
    \node at (1.5,1.5) {$d q\,d p$};
  \end{scope}

  \draw[green!45!black,->,line width=1.5pt] (5.25,2.1) -- (7.0,2.1)
    node[midway,above] {正则变换};

  % New coordinates: Q=q+0.55p, P=p.
  \begin{scope}[xshift=7.5cm]
    \draw[blue!70!black,->] (-.3,0) -- (6.5,0) node[right] {$Q$};
    \draw[blue!70!black,->] (0,-.3) -- (0,4.6) node[above] {$P$};
    \foreach \x in {1,2,3,4} {\draw[black!25] (\x,0) -- ({\x+2.2},4);}
    \foreach \y in {1,2,3,4} {\draw[black!25] ({.55*\y},\y) -- ({4+.55*\y},\y);}
    \fill[red!14] ({1+.55},1) -- ({2+.55},1) -- ({2+1.1},2) -- ({1+1.1},2) -- cycle;
    \draw[red] ({1+.55},1) -- ({2+.55},1) -- ({2+1.1},2) -- ({1+1.1},2) -- cycle;
    \node[red] at (2.35,1.5) {$dQ\,dP$};
    \node[blue!70!black] at (2.95,4.75) {新变量};
  \end{scope}

  \node[black!75,align=center] at (7.25,-1.0)
    {$dq\,dp=dQ\,dP$，相面积不变\\$[Q,P]=1$，正则结构不变};
\end{tikzpicture}
\end{document}
